\documentclass[11pt]{article}
\usepackage[margin=1in]{geometry}
\usepackage{amsmath,amssymb,amsthm,hyperref}
\hypersetup{colorlinks=true,urlcolor=blue}
\setlength{\emergencystretch}{2em}
\newtheorem{theorem}{Theorem}
\newtheorem{proposition}[theorem]{Proposition}
\title{Parameter concavity of the generalized inverse hyperbolic sine}
\author{Research record: MathDB 325512}
\date{12 September 2026}
\begin{document}
\maketitle

\section*{Task 1. Resolve the parameter-concavity assertion}
For $p>0$ and $x>0$, put
\[
 F_x(p)=\operatorname{arcsinh}_p(x)
       =\int_0^x(1+t^p)^{-1/p}\,dt.
\]
The MathDB assertion is that $F_x$ is strictly concave on $(0,\infty)$
for each $0<x<1$. This is also the domain in the published
Conjecture 2.4 of Baricz, Bhayo and Vuorinen \cite{BBV}.
A later survey \cite{survey}, Conjecture 5.1, states the narrower
domain $(1,\infty)$. We distinguish these two versions explicitly:
the published global assertion is false, and the narrower assertion
is true.

\begin{theorem}[An exact midpoint counterexample]\label{thm:counter}
At $x=1/4$, $p_-=1/32$, $p_0=1/16$, and $p_+=3/32$,
\[
 F_{1/4}(p_0)<\frac1{65536}
 <\frac1{16384}
 <\frac{F_{1/4}(p_-)+F_{1/4}(p_+)}2.
\]
Since $p_0=(p_-+p_+)/2$, this contradicts concavity.
\end{theorem}
\begin{proof}
The arithmetic--geometric mean inequality gives, for $t>0$,
\[
 1+t^p\ge2t^{p/2},\qquad
 (1+t^p)^{-1/p}\le 2^{-1/p}t^{-1/2}.
\]
Equality requires $t=1$, so integration gives
\begin{equation}\label{eq:flatbound}
 0<F_x(p)<2\sqrt{x}\,2^{-1/p}.
\end{equation}
For $x=1/4$ and $p=p_0$ this is $F_{1/4}(p_0)<2^{-16}$.
On $0<t<1/4$, we also have $1+t^{3/32}<2$. Consequently
\[
 (1+t^{3/32})^{-32/3}
 >2^{-32/3}>2^{-11},
\]
and therefore $F_{1/4}(p_+)>2^{-13}$.
Since $F_{1/4}(p_-)>0$, the midpoint average exceeds $2^{-14}$.
All bounds are strict and use only the stated rational parameters.
\end{proof}

\begin{proposition}[Failure for every fixed argument]
For every $x>0$, $F_x$ cannot be concave on any interval $(0,b)$
with $b>0$.
\end{proposition}
\begin{proof}
By \eqref{eq:flatbound}, $F_x(p)\to0$ and $F_x(p)/p\to0$ as
$p\downarrow0$. If $F_x$ were concave on $(0,b)$, then for
$0<r<s<b$, taking the limit of the concavity inequality with
one endpoint tending to zero would give
$F_x(r)\ge(r/s)F_x(s)$. Thus
$F_x(s)/s\le\lim_{r\downarrow0}F_x(r)/r=0$, contradicting positivity.
\end{proof}

\begin{theorem}[A positive range, including the survey conjecture]
\label{thm:positive}
For each $x>0$, $F_x$ is strictly increasing on $(0,\infty)$ and
strictly concave on $[\tfrac12\log2,\infty)$.
In particular, the conjecture restricted to $p>1$ is true.
\end{theorem}
\begin{proof}
For fixed $t>0$ set $h(p,t)=(1+t^p)^{-1/p}$,
$s=-p\log t$, and
\[
 A(s)=\log(1+e^{-s})+\frac{s}{1+e^s},
 \qquad B(s)=\frac{e^s}{(1+e^s)^2}.
\]
Direct differentiation gives
\[
 A'(s)=-sB(s),\qquad A(-s)=A(s),\qquad
 0<A(s)\le A(0)=\log2.
\]
The positivity and the limit $A(s)\to0$ as $|s|\to\infty$
also follow directly from the definition.
Since $\partial_p\log h=A(s)/p^2$, we have $\partial_p h>0$ and
\begin{equation}\label{eq:second}
 \frac{\partial_p^2h}{h}
 =\frac{A(s)^2-p\{2A(s)+s^2B(s)\}}{p^4}.
\end{equation}
If $p\ge\tfrac12\log2$, then
\[
 A(s)^2-2pA(s)\le A(s)\{A(s)-\log2\}\le0.
\]
Together with the nonpositive remaining term $-ps^2B(s)$,
this proves $\partial_p^2h\le0$. Equality can occur only when
$p=\tfrac12\log2$ and $t=1$. Integration over any interval
$(0,x)$ therefore gives $F_x''(p)<0$ throughout the stated range.

For completeness, differentiation under the integral is justified on
every compact parameter interval in $(0,\infty)$: $0<h\le1$,
$A$ is bounded, and $s^2B(s)$ is bounded on the real line.
The displayed first and second derivatives consequently have a
uniform integrable constant bound on the finite interval $(0,x)$.
\end{proof}

\paragraph{Scope.}
The counterexample settles the global ordinary-concavity question
actually listed in MathDB and the published paper. Theorem
\ref{thm:positive} also proves the narrower conjecture retained in
the later survey, so the result does not depend solely on a discrepancy
in the lower parameter endpoint. The survey's consequent
Tur\'an inequality for $p>2$ follows from strict concavity and the
arithmetic--geometric mean inequality.
The separate published Tur\'an assertion for every $p>1$ is not
refuted by the ordinary-concavity counterexample and is not needed
for the MathDB question.

\paragraph{Eligibility.}
The accompanying literature audit records fresh open/empty-claim
MathDB checks, the published source, the later survey, relevant
citing papers and public claim catalogues. No matching full public
claim was located in that bounded search.

\section*{Task 2. Determine the sharp concavity threshold}
For fixed $x>0$, determine the least $p_x$ such that $F_x$ is concave
on $[p_x,\infty)$. Find the best threshold uniform for $0<x<1$.

\section*{Task 3. Locate and count the inflection points}
Determine whether $F_x''$ has exactly one zero on $(0,\infty)$ for
each $x>0$, and describe how such zeros depend on $x$.

\section*{Task 4. Obtain uniform small-parameter asymptotics}
Derive a full asymptotic expansion of $2^{1/p}F_x(p)$ as $p\downarrow0$,
with explicit remainder bounds uniform on specified ranges of $x$,
including regimes in which $x$ also tends to zero.

\section*{Task 5. Analyze affine parameter paths}
For $\operatorname{arcsinh}_{p,q}(x)=
\int_0^x(1+t^q)^{-1/p}\,dt$, classify ordinary concavity along
$q=\alpha p+\beta$ on intervals where $p,q>0$, and identify sharp
conditions on $(\alpha,\beta)$ and $x$.

\begin{thebibliography}{9}
\bibitem{BBV}
\'A. Baricz, B. A. Bhayo and M. Vuorinen,
\emph{Tur\'an type inequalities for generalized inverse trigonometric functions},
Filomat \textbf{29}(2) (2015), 303--313.
\href{https://doi.org/10.2298/FIL1502303B}{doi:10.2298/FIL1502303B}.
Conjecture 2.4, p.~306.

\bibitem{survey}
L. Yin, L.-G. Huang, Y.-L. Wang and X.-L. Lin,
\emph{A survey for generalized trigonometric and hyperbolic functions},
Journal of Mathematical Inequalities \textbf{13}(3) (2019), 833--854.
\href{https://doi.org/10.7153/jmi-2019-13-58}{doi:10.7153/jmi-2019-13-58}.
Conjecture 5.1.
\end{thebibliography}
\end{document}
